\documentclass[exercice,a4]{thedocument}%
\usepackage{pstricks}%
\usepackage{array}%
\usepackage{textcomp}%
\startdatakeys
    \source{25GENMATNC1}
    \BO{2026}
    \cycle{4}
    \competencesDuSocle{CA,RA}
    \objectifsApprentissage{C4A1aa,C4A2ad,C4B2cb,C4D2aa,C4D3ac}
    \motsCles{droites parallèles,Thalès,aire,agrandissement,construction}
\enddatakeys
\begin{document}%
    
    \def\degres{\textsuperscript{\textdegree}}%

\medskip

La figure ci-dessous n'est pas à l'échelle.

\begin{minipage}{0.4\linewidth}
\begin{tabular}{|l|}\hline
On donne :\\
~\\ 
AE = 6 cm \\
AD = 15 cm\\
BE = 4 cm\\
$\widehat{\text{BEA}} = \widehat{\text{CDA}} = 110\degres$\\
Les points A, B et C sont alignés\\
Les points A, E et D sont alignés\\ \hline
\end{tabular}
\end{minipage}\hfill
\begin{minipage}{0.65\linewidth}
\psset{unit=0.72cm,arrowsize=2pt 3}
\begin{pspicture}(11.6,7)
%\psgrid
\pspolygon(0.2,1)(9,1)(11.45,6.7)%ADC
\psline(3.6,1)(4.4,3.16)%EB
\uput[l](0.2,1){A} \uput[ul](4.9,3.4){B} \uput[ur](11.5,6.7){C} \uput[dr](9,1){D} \uput[dr](3.6,1){E}
\psline[linewidth=0.6pt]{<->}(0.2,0.8)(3.6,0.8)\uput[d](1.9,0.9){\small 6 cm}
\psline[linewidth=0.6pt]{<->}(0.2,0.3)(9,0.3)\uput[d](4.6,0.3){\small 15 cm}
\psarc(9,1){0.6}{70}{180}\uput[ul](8.8,1.3){$110\degres$}
\psarc(3.6,1){0.6}{70}{180}\uput[ul](3.4,1.3){$110\degres$}
\rput{70}(4.25,2.08){4 cm}
\end{pspicture}
\end{minipage}

\medskip

\begin{enumerate}
\item Prouver que les droites(BE) et (CD)sont parallèles. 
\item Calculer CD.

\textbf{Rédiger la réponse en faisant apparaitre les différentes étapes.}
\item L'aire du triangle ABE, arrondie au dixième, est égale à $11,3$cm$^2$.
En déduire l'aire du triangle ACD arrondie au dixième.
\item Construire cette figure en vraie grandeur sur l'annexe 1.
\end{enumerate}

\bigskip
\showthesource
\end{document}
